%==============================================================
%  Chapter 1  --  Introduction
%==============================================================
\chapter{Introduction}
\label{ch:introduction}

Recommendation systems have become central to how users navigate large online
catalogues. On a marketplace with millions of items, no user can browse more than
a vanishing fraction of what is on offer, and the quality of the paths a platform
suggests --- ``users who viewed this also viewed'' and related item-to-item
surfaces --- largely determines how quickly a buyer reaches something relevant.
This thesis studies recommendation for IndiaMART, one of India's largest
business-to-business (B2B) marketplaces, where the task is to recommend products
given a product a user is currently looking at.

\section{Motivation and Setting}
\label{sec:intro-setting}

The IndiaMART setting brings together, in an acute form, the difficulties that
make industrial recommendation hard. The catalogue is enormous --- of the order
of tens of millions of live products spread across thousands of categories --- so
any method must scale to matrices and index structures with tens of millions of
rows. The interaction data is extremely \emph{sparse}: the vast majority of
product pairs are never seen together in a session, and the product-to-product
transition matrix has a density on the order of $10^{-4}$. The feedback is
entirely \emph{implicit} --- clicks, impressions, and downstream conversion
events collected from ordinary browsing, with no explicit ratings --- so the
signal is noisy and shaped by how items were presented. And the catalogue turns
over constantly, which keeps the \emph{cold-start} problem ever-present: newly
listed and long-tail products have little or no interaction history from which to
learn.

Against this backdrop, the incumbent production recommender is a simple but
strong frequency-based model (FBM): a product-to-product transition matrix built
from co-occurrence counts, row-normalised, and topped up through a
\emph{backfilling} step that fills empty recommendation slots with the platform's
pre-existing recommendations. This baseline is the yardstick used throughout the
thesis. It sets a deliberately high bar: a new method is only worth deploying if
it improves on the production system, not merely on some intermediate offline
metric. Much of what follows is a study of where richer models help against this
baseline, where they do not, and --- importantly --- why offline improvements do
not always translate into better recommendations.

\section{Approaches Studied}
\label{sec:intro-approaches}

The thesis develops three complementary lines of work.

\begin{itemize}
  \item \textbf{A $K$-step Markov-chain recommender via matrix factorization
        (Chapter~\ref{ch:markov-factorization}).} Product navigation is modelled
        as a Markov chain, and multi-hop recommendations are obtained from the
        powers of its transition matrix. Because those powers are infeasible to
        materialise at catalogue scale, they are expressed through a low-rank
        factorization, and several factorization schemes (ALS, NMF, and
        weighted NMF) are examined against the non-negativity and sparsity of the
        problem. The chapter reports both what the factorization saves and where
        it falls short of the frequency baseline.
  \item \textbf{Position-bias estimation and correction
        (Chapter~\ref{ch:position-bias}).} Because recommendations are consumed
        as ranked lists, raw clicks confound relevance with presentation: items
        shown near the top are clicked more simply for being shown near the top.
        This line of work models that distortion, estimates it from data through
        several independent routes, and folds the correction back into the
        recommendation engine as transition weights.
  \item \textbf{An embedding-based item-to-item model
        (Chapter~\ref{ch:embedding}).} Every product is given a dense embedding
        that is cold-started from its text metadata and then refined from observed
        interactions through a hierarchical cosine-softmax objective, so that
        content and collaborative signals are combined in a single vector space
        searched by approximate nearest neighbours. The chapter studies the
        model's training behaviour and, separately, its recommendation quality.
\end{itemize}

\section*{Key Contributions}
\addcontentsline{toc}{section}{Key Contributions}

The principal contributions of this thesis are the following.

\begin{enumerate}
  \item \textbf{A method for estimating and correcting position bias in the
        IndiaMART recommendation logs.} Adopting the examination hypothesis, the
        position bias is estimated by three independent routes --- a
        perfect-exposure aggregation, an anchoring of on-platform ranks to
        external traffic sources (Google and IndiaMART search), and a controlled
        A/B shifted-list experiment --- which agree qualitatively and thereby
        lend confidence to the estimate. The correction is integrated into the
        existing $K$-step recommender as per-transition-type weights, requiring no
        change to the engine itself. On held-out evaluation weeks this
        bias-weighted recommender improves click-through rate by roughly $3.5$ to
        $4$ percentage points consistently over the unweighted model.
  \item \textbf{An embedding-based item-to-item recommender that unifies content
        and collaborative signals at full catalogue scale.} A hierarchical
        cosine-softmax objective factorises the probability of a transition along
        the category taxonomy, making the softmax tractable over a catalogue of
        more than twenty million products; product embeddings are initialised from
        text metadata and refined from interactions; and the learned space is
        served through approximate nearest-neighbour search. Beyond the model
        itself, the chapter contributes a careful evaluation showing that the
        interaction training adds real signal over a content-only baseline, and a
        cautionary finding that the training objective can move opposite to
        recommendation quality.
  \item \textbf{An analysis of matrix-factorization approaches to the $K$-step
        Markov recommender, and their limitations.} The factorized formulation
        yields large memory savings for higher-order transitions, but an
        empirical evaluation shows that it does not preserve recommendation
        quality relative to the raw transition matrix or the production baseline,
        and a structural limitation of weighted NMF --- unconstrained,
        unobserved entries that inflate and dominate the Markov chain --- is
        identified.
  \item \textbf{A cross-cutting, practically important observation:} offline
        optimisation metrics --- reconstruction error for the factorization,
        negative log-likelihood for the embedding model --- are unreliable
        proxies for recommendation quality in this setting, and the simple
        frequency-based model with backfilling remains a strong baseline that the
        richer models studied here do not uniformly surpass.
\end{enumerate}

\section{Organisation of the Thesis}
\label{sec:intro-organization}

The remainder of the thesis is organised as follows.
Chapter~\ref{ch:datasets} describes the datasets --- the click-stream logs, the
transition data derived from them, and the product catalogue --- together with
the preprocessing common to all later chapters.
Chapter~\ref{ch:markov-factorization} develops the $K$-step Markov-chain
recommender and its matrix-factorization formulation, and evaluates the
factorization against the frequency baseline.
Chapter~\ref{ch:position-bias} takes up position bias: its modelling,
estimation, and correction, and the integration of that correction into the
recommendation engine.
Chapter~\ref{ch:embedding} presents the embedding-based item-to-item model, its
implementation at scale, and its recommendation-quality evaluation.
Chapter~\ref{ch:conclusion} concludes, drawing together the findings and
outlining directions for further work.
